\documentclass[../../../include/open-logic-section]{subfiles}

\begin{document}

\olfileid{sth}{z}{sep}
\olsection{પૃથક્કરણ}

સહજ ગણરચનાને બદલવા માટેના સિદ્ધાંતથી શરૂઆત કરીએ:

\begin{axiom}[પૃથક્કરણનું પ્રરૂપ]
દરેક સૂત્ર $\phi(x)$ માટે આ એક સ્વયંસિદ્ધિ છે: કોઈ પણ $A$ માટે
ગણ $\Setabs{x \in A}{\phi(x)}$ અસ્તિત્વ ધરાવે છે.
\end{axiom}

નોંધો કે આ એક સ્વયંસિદ્ધિ નથી. તે સ્વયંસિદ્ધિઓનું \emph{પ્રરૂપ}
છે. પૃથક્કરણની \emph{અનંત સંખ્યામાં} સ્વયંસિદ્ધિઓ છે: દરેક સૂત્ર
$\phi(x)$ માટે એક. આ પ્રરૂપને નીચે પ્રમાણે પણ લખી શકાય
(અને સામાન્ય રીતે એમ જ લખાય છે):

\begin{defish}
જેમાં ``$S$'' ન આવે એવા કોઈ પણ સૂત્ર $\phi(x)$ માટે આ એક
સ્વયંસિદ્ધિ છે:
\[
\forall A \exists S \forall x(x \in S \liff (\phi(x) \land x \in A)).
\]
\end{defish}

\olref[sth][][]{part}ની શરૂઆતમાં જણાવેલી રીત મુજબ પૃથક્કરણની
સ્વયંસિદ્ધિઓનાં સૂત્રો~$\phi$ પ્રાચલો ધરાવી શકે છે.
\footnote{આનો અર્થ સમજવા માટે
\olref[sfr][infinite][induction]{natinductionschema} પછીની તરતની
ચર્ચા જુઓ.}

ગણોની જે સંચયી-પુનરાવર્તી કલ્પનાનું વર્ણન કર્યું છે તે તરત જ
પૃથક્કરણને યોગ્ય ઠરાવે છે. તેનું કારણ જોવા માટે $A$ને ગણ લો.
તેથી $A$ કોઈ તબક્કા~$S$ પર રચાય છે (\stageshier પ્રમાણે).
$A$ તબક્કા~$S$ પર રચાયો હોવાથી $A$ના બધા સભ્યો તબક્કા $S$
પહેલાં રચાયા હતા (\stagesacc પ્રમાણે). હવે ખાસ કરીને એવા બધા
ગણોનો વિચાર કરો જે $A$ના સભ્યો છે અને $\phi$ને પણ સંતોષે છે;
સ્પષ્ટ છે કે આ બધા ગણો પણ તબક્કા~$S$ પહેલાં રચાયા હતા.
તેથી તબક્કા~$S$ પર તેઓનો ગણ $\Setabs{x \in A}{\phi(x)}$
પણ રચાય છે (\stagesacc પ્રમાણે).

સહજ ગણરચનાથી વિપરીત, આ રસેલનો વિરોધાભાસ ટાળે છે. કારણ કે
ગણ $\Setabs{x}{x \notin
x}$ના અસ્તિત્વનો સીધો દાવો કરી શકતા નથી. તેના બદલે કોઈ
ગણ~$A$ \emph{આપેલો હોય}, તો ગણ
$R_A = \Setabs{x \in A}{x \notin x}$ના અસ્તિત્વનો દાવો કરી
શકીએ. પરંતુ આમાંથી માત્ર $R_A \notin R_A$ અને $R_A \notin A$
જ સાબિત થાય છે, અને તેમાંથી કોઈ વાત ખાસ ચિંતાજનક નથી.

જોકે પૃથક્કરણનું એક તરત મળતું અને નોંધપાત્ર પરિણામ છે:

\begin{thm}\ollabel{thm:NoUniversalSet}
કોઈ \emph{સાર્વત્રિક} ગણ નથી; એટલે કે $\Setabs{x}{x = x}$
અસ્તિત્વ ધરાવતો નથી.
\end{thm}

\begin{proof}
પરસ્પરવિરોધ મેળવવા માટે ધારો કે $V$ સાર્વત્રિક ગણ છે. તો
પૃથક્કરણથી $R =
\Setabs{x \in V}{x \notin x} = \Setabs{x}{x \notin x}$ અસ્તિત્વ
ધરાવે છે, જે રસેલના વિરોધાભાસની વિરુદ્ધ છે.
\end{proof}

સાર્વત્રિક ગણ ન હોવો---ખરેખર, ગણોના તબક્કાવાર સ્તરક્રમને કોઈ
અંતિમ સીમા ન હોવી---સંચયી-પુનરાવર્તી કલ્પનાના સૌથી મૂળભૂત
વિચારોમાંનો એક છે. તેથી સહજ રીતે જુદા માર્ગે પણ ત્યાં પહોંચી
શકીએ છીએ તે જોવું યોગ્ય છે. સાર્વત્રિક ગણ પોતાનો જ
!!a{element} હોવો જોઈએ. પરંતુ આપણી સંચયી-પુનરાવર્તી કલ્પના મુજબ
દરેક ગણ તેના બધા !!{element}s પછીના સૌપ્રથમ તબક્કે સ્તરક્રમમાં
(પ્રથમ વાર) દેખાય છે. આથી \emph{કોઈ} ગણ સ્વસભ્ય નથી એ ફલિત થાય
છે. કારણ કે કોઈ સ્વસભ્ય ગણ પહેલી વાર દેખાયો હોય તે તબક્કા
પછીના સૌપ્રથમ તબક્કે જ તેને પહેલી વાર દેખાવું પડે, જે અસંગત છે.
(\olref[sth][spine][rank]{defnsetrank}માં ગણના ``સ્તરાંક''ના
ખ્યાલ વડે આ સમજૂતીને વધુ ચોક્કસ કેવી રીતે બનાવવી તે જોઈશું.
પરંતુ તે માટે પહેલાં થોડી વધુ સ્વયંસિદ્ધિઓ જોઈએ.)

પૃથક્કરણ અને ઘટકો દ્વારા નિર્ધારિત સમાનતાનાં થોડાં વધુ પરિણામો
આ છે.

\begin{prop}\ollabel{prop:emptyexists}
જો કોઈ ગણ અસ્તિત્વ ધરાવતો હોય, તો $\emptyset$ અસ્તિત્વ ધરાવે છે.
\end{prop}

\begin{proof}
જો $A$ ગણ હોય, તો પૃથક્કરણથી
$\emptyset = \Setabs{x \in A}{x \neq x}$ અસ્તિત્વ ધરાવે છે.
\end{proof}

\begin{prop}
કોઈ પણ ગણો $A$ અને $B$ માટે $A \setminus B$ અસ્તિત્વ ધરાવે છે.
\end{prop}

\begin{proof}
પૃથક્કરણથી $A \setminus B = \Setabs{x \in A}{x \notin B}$
અસ્તિત્વ ધરાવે છે.
\end{proof}

(લગભગ) મનસ્વી છેદગણો પણ અસ્તિત્વ ધરાવે છે:

\begin{prop}\ollabel{prop:intersectionsexist}
જો $A \neq \emptyset$, તો
$\bigcap A = \Setabs{x}{(\forall y \in A)x \in y}$ અસ્તિત્વ ધરાવે છે.
\end{prop}

\begin{proof}
$A \neq \emptyset$ લો; તેથી કોઈ $c \in A$ છે. તો
$\bigcap A =
\Setabs{x}{(\forall y \in A)x \in y} = \Setabs{x \in c}{(\forall y \in
A)x \in y}$, જે પૃથક્કરણથી અસ્તિત્વ ધરાવે છે.
\end{proof}

જોકે $A \neq \emptyset$ હોવાની શરત ધ્યાનમાં રાખો; કારણ કે
$\bigcap
\emptyset$ રિક્તતાને કારણે સાર્વત્રિક ગણ હોત, જે
\olref{thm:NoUniversalSet}ની વિરુદ્ધ છે.

\end{document}
